\documentclass{stex}

\libusepackage{naproche}
\libinput{preamble}

\begin{document}
\begin{smodule}{successor-ordinals.ftl}

\importmodule[libraries/set-theory]{ordinals.ftl}

\symdef{Succ}{\textrm{succ}}
\symdef{Pred}{\textrm{pred}}

\symdecl*{successor ordinal}

\begin{definition}[forthel,for=Succ]
  Let $\alpha$ be an ordinal.
  $\emph{\Succ(\alpha)} \DefEq \alpha \Union \Singleton{\alpha}$.
\end{definition}

\begin{proposition}[forthel]
  Let $\alpha$ be an ordinal.
  Then $\Succ(\alpha)$ is an ordinal.
\end{proposition}
\begin{proof}[forthel]
  (1) $\Succ(\alpha)$ is transitive.
  \begin{proof}
    Let $x \Elem \Succ(\alpha)$ and $y \Elem x$.
    Then $x \Elem \alpha$ or $x \Eq \alpha$.
    Hence $y \Elem \alpha$.
    Thus $y \Elem \Succ(\alpha)$.
  \end{proof}

  (2) Every element of $\Succ(\alpha)$ is transitive.
  \begin{proof}
    Let $x \Elem \Succ(\alpha)$.
    Then $x \Elem \alpha$ or $x \Eq \alpha$.
    Hence $x$ is transitive.
    Indeed $\alpha$ is transitive and every element of $\alpha$ is transitive.
  \end{proof}
\end{proof}

\begin{definition}[forthel,for=successor ordinal]
  A \emph{successor ordinal} is an ordinal $\alpha$ such that $\alpha \Eq \Succ(\beta)$ for some ordinal $\beta$.
\end{definition}

\begin{proposition}[forthel]
  Let $\alpha, \beta$ be ordinals.
  If $\Succ(\alpha) \Eq \Succ(\beta)$ then $\alpha \Eq \beta$.
\end{proposition}
\begin{proof}[forthel]
  Assume $\Succ(\alpha) \Eq \Succ(\beta)$.

  (1) $\alpha \Subclass \beta$.
  \begin{proof}
    Let $\gamma \Elem \alpha$.
    Then $\gamma \Elem \alpha \Union \Singleton{\alpha}
      \Eq \Succ(\alpha)
      \Eq \Succ(\beta)
      \Eq \beta \Union \Singleton{\beta}$.
    Hence $\gamma \Elem \beta$ or $\gamma \Eq \beta$.
    Assume $\gamma \Eq \beta$.
    Then $\beta \Elem \alpha$.
    Hence $\beta
      \Eq (\beta \Union \Singleton{\beta}) \Diff \Singleton{\gamma}
      \Eq (\alpha \Union \Singleton{\alpha}) \Diff \Singleton{\gamma}
      \Eq (\alpha \Diff \Singleton{\gamma}) \Union \Singleton{\alpha}$.
    Therefore $\alpha \Elem \beta$.
    Consequently $\alpha \Elem \beta \Elem \alpha$.
    Contradiction.
  \end{proof}

  (2) $\beta \Subclass \alpha$.
  \begin{proof}
    Let $\gamma \Elem \beta$.
    Then $\gamma \Elem \beta \Union \Singleton{\beta}
      \Eq \Succ(\beta)
      \Eq \Succ(\alpha)
      \Eq \alpha \Union \Singleton{\alpha}$.
    Hence $\gamma \Elem \alpha$ or $\gamma \Eq \alpha$.
    Assume $\gamma \Eq \alpha$.
    Then $\alpha \Elem \beta$.
    Hence $\alpha
      \Eq (\alpha \Union \Singleton{\alpha}) \Diff \Singleton{\gamma}
      \Eq (\beta \Union \Singleton{\beta}) \Diff \Singleton{\gamma}
      \Eq (\beta \Diff \Singleton{\gamma}) \Union \Singleton{\beta}$.
    Therefore $\beta \Elem \alpha$.
    Consequently $\beta \Elem \alpha \Elem \beta$.
    Contradiction.
  \end{proof}
\end{proof}

\begin{definition}[forthel,for=Pred]
  Let $\alpha$ be a successor ordinal.
  $\emph{\Pred(\alpha)}$ is the ordinal $\beta$ such that $\alpha \Eq \Succ(\beta)$.
\end{definition}

\end{smodule}
\end{document}
